% from ega0-14-fr.tex
\oldpage[0\textsubscript{IV-1}]{6}
\begin{env}[14.1.1]
\label{0.14.1.1-fr}
Soit $I$ un ensemble ordonné~; une \emph{chaîne} d'éléments de $I$ est par
définition une suite finie strictement croissante
$i_0<i_1<\cdots<i_n$ d'éléments de $I$ ($n\geq 0$)~; par définition la
\emph{longueur} de cette chaîne est $n$. Si $X$ est un espace topologique,
l'ensemble de ses parties \emph{fermées irréductibles} est ordonné par
inclusion, d'où la notion de \emph{chaîne} de parties fermées irréductibles de
$X$.
\end{env}

\begin{definition}[14.1.2]
\label{0.14.1.2-fr}
Soit $X$ un espace topologique. On appelle \emph{dimension combinatoire de
$X$} (ou simplement \emph{dimension de $X$} s'il ne peut y avoir de
confusion) et on note $\dim_c(X)$, ou simplement $\dim(X)$, la borne
supérieure des longueurs des chaînes de parties fermées irréductibles de $X$.
Pour tout $x\in X$, on appelle \emph{dimension combinatoire de $X$ en $x$}
(ou simplement \emph{dimension de $X$ en $x$}) et on note $\dim_x(X)$ le
nombre $\inf_U(\dim(U))$, où $U$ parcourt l'ensemble des voisinages ouverts de
$x$ dans $X$.
\end{definition}

Il résulte de cette définition que l'on a
\[
  \dim(\varnothing)=-\infty
\]
(la borne supérieure dans $\overline{\mathbb R}$ de la partie vide étant
$-\infty$). Si $(X_\alpha)$ est la famille des composantes irréductibles de
$X$, on a
\[
  \dim(X)=\sup_\alpha\dim(X_\alpha)
  \tag{14.1.2.1}\label{0.14.1.2.1-fr}
\]
car toute chaîne de parties fermées irréductibles de $X$ est par définition
contenue dans une composante irréductible de $X$, et inversement ces
composantes sont fermées dans $X$, donc toute partie fermée irréductible d'un
$X_\alpha$ est une partie fermée irréductible de $X$.

\begin{definition}[14.1.3]
\label{0.14.1.3-fr}
On dit qu'un espace topologique $X$ est \emph{équidimensionnel} si toutes ses
composantes irréductibles ont même dimension (égale alors à $\dim(X)$ par
(14.1.2.1)).
\end{definition}

\begin{proposition}[14.1.4]
\label{0.14.1.4-fr}
\begin{enumerate}
  \item[\rm(i)] Pour toute partie $Y$ d'un espace topologique $X$, on a
  $\dim(Y)\leq\dim(X)$.
  \item[\rm(ii)] Si un espace topologique $X$ est réunion finie de parties
  fermées $X_i$, on a $\dim(X)=\sup_i\dim(X_i)$.
\end{enumerate}
\end{proposition}

\begin{proof}
Pour toute partie fermée irréductible $Z$ de $Y$, l'adhérence $\overline Z$ de
$Z$ dans $X$ est irréductible ($0_{\mathrm I}$, 2.1.2) et
$\overline Z\cap Y=Z$, d'où (i). D'autre part, si
$X=\bigcup_{i=1}^n X_i$, où les $X_i$ sont fermés, toute partie fermée
irréductible de $X$ est contenue dans un des $X_i$ ($0_{\mathrm I}$, 2.1.1),
et par suite toute chaîne de parties fermées irréductibles de $X$ est contenue
dans un des $X_i$~; d'où (ii).
\end{proof}

On déduit de (14.1.4, (i)) que pour tout $x\in X$, on peut aussi écrire
\[
  \dim_x(X)=\lim_U\dim(U)
  \tag{14.1.4.1}\label{0.14.1.4.1-fr}
\]
la limite étant prise suivant l'ensemble filtrant décroissant des voisinages
ouverts de $x$ dans $X$.

